\scanpage{115}
induit sur $H$ et $H'$ l'automorphisme
\[
x\longmapsto x'=\operatorname{Tr}(x)\cdot1-\unsure{x|a}
\]

En tous cas, les isom qui envoient $(H,H')$ sur $(\overline H,\overline H')$ forment un torseur sous le groupe $\mathbb Z/2\mathbb Z$ (si $\rho_0+4$ inv.) lequel torseur peut (je présume) être non trivial (même sachant que le torseur associé de groupe $\underline{\operatorname{Aut}}(A)\simeq\deleted{\unsure{A^*/}}\check V(A)^*/G_m$ est trivial

\deleted{\unsure{On|En}} En effet : [Si $P$ est ce \inserted{$\mathbb Z/2$-}torseur sur $\operatorname{Spec} k$, on aura un hom
\[
\inserted{P\xrightarrow{\sim}\operatorname{isom}(\unsure{\xi|\mathcal E|\varepsilon}_H,\xi_{\overline H})\;(\simeq\xi_H\wedge\xi_{\overline H})}
\]
\[
\deleted{P\times\xi_H\longrightarrow\xi_{\overline H}}
\]
% slanted
\margin{$\xi_H$ désigne le rev. quadratique étale de $S=\operatorname{Spec}k$ défini par $H$} \deleted{qui définit} et de même
\[
P\xrightarrow{\sim}\xi_{H'}\wedge\xi_{\overline H'}\quad]
\]

\deleted{\unsure{Donc}} $P$ n'a de sections que si $\xi_{H\deleted{\unsure{'}}}\simeq\xi_{\overline H}$, \inserted{$\xi_{H'}\simeq\xi_{\overline{H'}}$} ; Prenant par exemple $H'=\overline{H'}$ correspondant à la bisection $\{0,\infty\}$ de $\mathbb P^1$, et $H,\overline H$ étant les bisections définies resp.\textsuperscript{t} par
% slanted
\margin{sur un corps\ldots}
\[
\begin{aligned}
z^2-bz+c&=0 & b^2-4c&\mathrel{\deleted{\unsure{=}}\ne}0\\ z^2-\overline b z+\overline c&=0 & \overline b^2-4\overline c&\ne0
\end{aligned}
\]
% archive sheet number: 114
